\documentclass[a4paper,12pt]{article}
\usepackage[T2A]{fontenc}
\usepackage[utf8]{inputenc}
\usepackage[russian]{babel}
\usepackage{amsmath,amssymb,mathtools}
\usepackage{geometry}
\usepackage{microtype}
\usepackage{tikz}
\usetikzlibrary{arrows.meta,positioning,shapes.geometric,calc}
\usepackage{pgfplots}
\pgfplotsset{compat=1.18}
\usepackage{booktabs,tabularx,array}
\usepackage{enumitem}
\usepackage{fancyhdr}
\usepackage{hyperref}
\usepackage{parskip}
\usepackage{setspace}
\usepackage{xcolor}
\usepackage{needspace}
\usepackage{icomma}
\geometry{left=22mm,right=22mm,top=24mm,bottom=24mm}
\definecolor{accent}{HTML}{008080}
\definecolor{darkaccent}{HTML}{006469}
\definecolor{textgray}{HTML}{3C4345}
\definecolor{pale}{HTML}{91B9BA}
\color{textgray}
\setstretch{1.12}
\setlength{\parindent}{0pt}
\setlength{\parskip}{4pt}
\setlist[itemize]{leftmargin=5mm,itemsep=2.5pt,topsep=3pt}
\setlist[enumerate]{leftmargin=7mm,itemsep=4pt,topsep=4pt}
\hypersetup{colorlinks=true,linkcolor=darkaccent,urlcolor=darkaccent,pdfauthor={Лёвин Артём Александрович},pdftitle={Чек-лист. Кусочные функции, уравнения и задачи на движение}}
\pagestyle{fancy}
\fancyhf{}
\fancyhead[L]{\footnotesize Кусочные функции и уравнения}
\fancyhead[R]{\footnotesize София Кальней}
\renewcommand{\headrulewidth}{0.35pt}
\renewcommand{\headrule}{\hbox to\headwidth{\color{accent}\leaders\hrule height \headrulewidth\hfill}}
\fancyfoot[C]{\small\thepage}
\setlength{\headheight}{24pt}
\newcommand{\secname}[1]{\par\vspace{5pt}{\color{accent}\Large #1}\par\vspace{3pt}\noindent{\color{accent}\rule{\linewidth}{0.45pt}}\vspace{2pt}}
\newcommand{\keyidea}[1]{{\color{darkaccent}\textbf{Главная идея.}} #1}
\newcommand{\ex}[1]{\needspace{7\baselineskip}\vspace{2pt}\textbf{#1}\quad}
\newcommand{\task}[2]{\needspace{6\baselineskip}\noindent\textbf{#1.}\enspace #2\par\vspace{9pt}}
\begin{document}
\begin{titlepage}
\thispagestyle{empty}
\begin{center}
\vspace*{25mm}
{\color{darkaccent}\large МАТЕМАТИКА}\par
\vspace{16mm}
{\LARGE Чек-лист}\par
\vspace{11mm}
{\color{accent}\LARGE Кусочные функции,\par пересечения графиков\par и уравнения}\par
\vspace{12mm}
{\large Подготовка к ОГЭ по математике}\par
\vfill
{\large София Кальней}\par
\vspace{12mm}
{1 ноября 2026 г.}\par
\vspace{16mm}
{\small Лёвин Артём Александрович}\par
{\small эксклюзивно для Софии Кальней}\par
\vspace{12mm}
\begin{tikzpicture}
\draw[accent,line width=0.8pt] (0,0) .. controls (1.3,0.68) and (2.9,-0.65) .. (4.4,0)
 .. controls (5.6,0.53) and (7.0,-0.50) .. (8.6,0);
\end{tikzpicture}
\vspace{5mm}
\end{center}
\end{titlepage}
\setcounter{page}{1}

\secname{1. Замена переменной в уравнении}

Когда в уравнении несколько раз повторяется выражение \(\dfrac1x\), удобно заменить его одной буквой:
\[
 t=\frac1x,\qquad \frac1{x^2}=t^2,\qquad x\ne0.
\]
\keyidea{Сначала решите уравнение относительно новой переменной, затем выполните обратную замену.}

\ex{Пример.} Решим уравнение
\[
\frac1{x^2}-\frac1x-6=0,\qquad x\ne0.
\]
Положим \(t=\dfrac1x\). Тогда
\[
 t^2-t-6=0,\qquad (t-3)(t+2)=0.
\]
Получаем \(t=3\) или \(t=-2\). Возвращаемся к \(x\):
\[
 \frac1x=3\Rightarrow x=\frac13,\qquad
 \frac1x=-2\Rightarrow x=-\frac12.
\]
\textbf{Ответ:} \(\displaystyle -\frac12;\ \frac13\).

\textcolor{darkaccent}{Проверьте:} \(x\ne0\); оба найденных корня допустимы. Если заменено \(t=1/x\), то при обратной замене \(x=1/t\), а не \(x=t\).

\secname{2. Движение навстречу и остановка}

Для равномерного движения используйте \(S=vt\). Если участники выехали одновременно и один из них остановился на \(\tau\) часа, время \emph{движения} остановившегося будет на \(\tau\) меньше.

Пусть первый велосипедист ехал со скоростью \(v_1\) и сделал остановку, второй двигался без остановок со скоростью \(v_2\). Если встретились через \(t\) часов после выезда, то
\[
\underbrace{v_1(t-\tau)}_{\text{первый}}+
\underbrace{v_2t}_{\text{второй}}=S,\qquad t\ge\tau.
\]

\begin{center}
\renewcommand{\arraystretch}{1.22}
\begin{tabular}{@{}lccc@{}}
\toprule
Участник & Скорость & Время движения & Путь\\
\midrule
Первый & \(v_1\) & \(t-\tau\) & \(v_1(t-\tau)\)\\
Второй & \(v_2\) & \(t\) & \(v_2t\)\\
\bottomrule
\end{tabular}
\end{center}

\textbf{Важно:} переводите минуты в часы; после нахождения \(t\) вычисляйте именно ту величину, которую спрашивают в условии.

\clearpage
\secname{3. Кусочная функция: каждому фрагменту своя область}

Кусочная функция задаётся разными формулами на разных промежутках. Стройте \textbf{каждый фрагмент отдельно} и оставляйте только ту часть графика, которая разрешена условием.

\ex{Пример.}
\[
 f(x)=\begin{cases}
 -\dfrac6x,&x<-2\ \text{или}\ x>1,\\[4pt]
 -3x-3,&-2\le x\le1.
 \end{cases}
\]

На внешних промежутках получаем части гиперболы \(y=-6/x\), на среднем --- отрезок прямой \(y=-3x-3\).

\begin{center}
\begin{tikzpicture}
\begin{axis}[
 width=12.4cm,height=7.0cm,
 xmin=-5.4,xmax=5.4,ymin=-7.4,ymax=6.0,
 axis lines=middle,axis line style={black!70},
 axis equal image=false,
 xtick={-4,-2,0,1,2,4},ytick={-6,-3,0,3,6},
 ticklabel style={font=\scriptsize},xlabel={$x$},ylabel={$y$},
 xlabel style={at={(axis description cs:0.99,0.55)},anchor=south east},
 ylabel style={at={(axis description cs:0.50,0.99)},anchor=north west},
 grid=both,minor tick num=0,grid style={draw=black!8},
 clip=true]
\addplot[accent,thick,samples=150,domain=-5.4:-2.001] {-6/x};
\addplot[accent,thick,samples=150,domain=1.001:5.4] {-6/x};
\addplot[darkaccent,very thick,domain=-2:1] {-3*x-3};
\addplot[only marks,mark=*,mark size=2.2pt,color=darkaccent] coordinates {(-2,3) (1,-6)};
\end{axis}
\end{tikzpicture}
\end{center}

\textbf{Граничные точки:} при \(x=-2\) обе формулы дают \(y=3\), а при \(x=1\) --- \(y=-6\). График соединяется в этих точках: обе точки принадлежат среднему фрагменту.

\secname{4. Закрашенные и выколотые точки}

\begin{itemize}
\item Если граница входит в область фрагмента (\(\le\), \(\ge\)), точка \textbf{закрашенная}.
\item Если граница исключена (\(<\), \(>\)), точка \textbf{выколотая} для этого фрагмента.
\item Если \(y=a/x\), значение \(x=0\) запрещено при любом \(a\ne0\).
\item В одной граничной абсциссе разные формулы могут давать \textbf{разные ординаты}. Тогда одна точка может быть выколотой, а другая --- закрашенной.
\end{itemize}

\clearpage
\secname{5. Прямая \texorpdfstring{\(y=kx\)}{y=kx} и число пересечений}

Прямая \(y=kx\) всегда проходит через начало координат. При изменении \(k\) меняется её наклон: \(k>0\) --- прямая возрастает, \(k<0\) --- убывает, \(k=0\) --- ось \(Ox\).

\keyidea{Число общих точек --- число различных решений уравнения \(f(x)=kx\) с обязательной проверкой промежутков.}

\textbf{Для гиперболы} \(y=a/x\) (\(a\ne0\)):
\[
\frac ax=kx\quad\Longrightarrow\quad kx^2=a,\qquad x\ne0.
\]
При \(k\ne0\) возможны корни \(x=\pm\sqrt{a/k}\), только если \(a/k>0\). Каждый корень должен входить в разрешённый промежуток гиперболы.

\textbf{Для прямой} \(y=mx+b\):
\[
 mx+b=kx\quad\Longrightarrow\quad(m-k)x=-b.
\]
Если \(k=m\) и \(b\ne0\), прямые параллельны и общих точек нет. Если \(k=m\) и \(b=0\), линии совпадают на разрешённом фрагменте, поэтому общих точек бесконечно много, если фрагмент содержит промежуток.

\textbf{Граничные значения параметра.} Если прямая \(y=kx\) проходит через точку \((x_0,y_0)\), где \(x_0\ne0\), то
\[
 k=\frac{y_0}{x_0}.
\]
В таком положении число общих точек часто меняется. Проверяйте строгость границ отдельно.

\ex{Пример: ровно три точки.} Для функции со страницы 2 найдём \(k\), при котором \(y=kx\) имеет с её графиком ровно три общие точки.

Из \(-6/x=kx\) получаем \(x^2=-6/k\); две допустимые точки на внешних ветвях существуют, если
\[
 -\frac32<k<0.
\]
На этом интервале пересечение со средним отрезком имеет абсциссу
\[
 -3x-3=kx\quad\Longrightarrow\quad x=-\frac{3}{k+3}\in(-2,-1).
\]
Значит, добавляется третья точка. На границе \(k=-3/2\) остаются две различные точки: одна из них лежит на месте соединения фрагментов.

\textbf{Ответ:} \(\displaystyle k\in\left(-\frac32;0\right)\).

\clearpage
\secname{6. Если кусочная функция имеет разрывы}

Рассмотрим функцию
\[
 g(x)=\begin{cases}
 -2x+1,&x\le-2\ \text{или}\ x\ge2,\\[4pt]
 \dfrac2x,&-2<x<0\ \text{или}\ 0<x<2.
\end{cases}
\]
Внутренняя гипербола имеет выколотые точки \((-2,-1)\), \((2,1)\). У внешней прямой закрашены точки \((-2,5)\), \((2,-3)\).

\textbf{Шаг 1.} Пересечения с гиперболой:
\[
\frac2x=kx\quad\Longrightarrow\quad x^2=\frac2k.
\]
Оба решения попадают внутрь \((-2;2)\), не включая \(0\), только при \(k>\frac12\). При \(k=\frac12\) получаются \(x=\pm2\), которые для гиперболы исключены.

\textbf{Шаг 2.} Пересечения с внешней прямой:
\[
-2x+1=kx\quad\Longrightarrow\quad x=\frac1{k+2},\qquad k\ne-2.
\]
Условие \(x\le-2\) или \(x\ge2\) выполняется ровно при
\[
 k\in\left[-\frac52;-2\right)\cup\left(-2;-\frac32\right].
\]
При \(k=-2\) прямые параллельны и пересечений с внешним фрагментом нет.

\textbf{Шаг 3.} Собираем те значения, при которых \emph{ни один} фрагмент не дал общей точки:
\[
\boxed{\displaystyle k\in\left(-\infty;-\frac52\right)\cup\{-2\}\cup
\left(-\frac32;\frac12\right].}
\]

\textcolor{darkaccent}{Здесь особенно важно} различать выколотые и закрашенные точки и отдельно проверять случай параллельности прямых.

\secname{7. Типичные ошибки}
\begin{itemize}[itemsep=0pt]
\item Не проверяют, попадает ли найденный корень в промежуток своего фрагмента.
\item Учитывают выколотую точку или дважды считают одну граничную точку.
\item Пропускают случай параллельных или совпадающих прямых.
\item Путают \(x\) и \(1/x\) при обратной замене либо время движения и время от выезда.
\end{itemize}

\clearpage
\begin{center}
{\color{accent}\LARGE Алгоритм: сколько общих точек?}\par
\vspace{7mm}
\end{center}
\begin{center}
\begin{tikzpicture}[
 node distance=9mm,
 every node/.style={font=\normalsize,align=center},
 action/.style={draw=darkaccent,rounded corners=2.5pt,text width=11.3cm,minimum height=12mm,inner sep=5pt,line width=0.65pt},
 terminal/.style={draw=accent,ellipse,text width=8.8cm,minimum height=10mm,inner sep=5pt,line width=0.75pt},
 link/.style={-{Latex[length=2.9mm]},line width=0.85pt,draw=darkaccent}
 ]
\node[terminal] (start) {Даны график кусочной функции \(y=f(x)\) и прямая \(y=kx\)};
\node[action,below=of start] (split) {Разделить функцию на фрагменты и записать область каждого фрагмента};
\node[action,below=of split] (eq) {Для каждого фрагмента отдельно решить уравнение \(f_i(x)=kx\)};
\node[action,below=of eq] (check) {Оставить только корни, входящие в область соответствующего фрагмента};
\node[action,below=of check] (ends) {Проверить граничные точки, выколы, параллельность и совпадение прямых};
\node[action,below=of ends] (count) {Найти критические значения \(k\) и подсчитать \textbf{различные} общие точки на каждом промежутке};
\node[terminal,below=of count] (done) {Выбрать только значения \(k\), дающие требуемое число пересечений};
\draw[link] (start)--(split);
\draw[link] (split)--(eq);
\draw[link] (eq)--(check);
\draw[link] (check)--(ends);
\draw[link] (ends)--(count);
\draw[link] (count)--(done);
\end{tikzpicture}
\end{center}
\vspace{5mm}
\begin{center}
\textcolor{darkaccent}{Контроль: каждый найденный корень должен удовлетворять и формуле, и условию на \(x\).}
\end{center}

\clearpage
\secname{Тренировочный блок --- Путь А}
\textit{10 заданий: от повторения основных приёмов к исследованию параметра.}

\textbf{Базовый уровень}
\task{1}{Решите уравнение: \(\displaystyle \frac1{x^2}-\frac5x+6=0\).}
\task{2}{Решите уравнение заменой переменной: \(\displaystyle \frac9{x^2}+\frac3x-2=0\).}

\textbf{Задачи на движение}
\task{3}{Из двух городов, расстояние между которыми 114 км, одновременно навстречу друг другу выехали два велосипедиста со скоростями 12 км/ч и 18 км/ч. Первый сделал остановку на 30 минут и успел продолжить движение до встречи. Сколько километров проехал \textit{второй} велосипедист до встречи?}
\task{4}{Из пунктов, расстояние между которыми 132 км, одновременно навстречу друг другу выехали два велосипедиста со скоростями 15 км/ч и 20 км/ч. Второй сделал остановку на 24 минуты и затем продолжил движение. Сколько километров проехал \textit{первый} велосипедист до встречи?}

\textbf{Кусочные функции}
\task{5}{Постройте график функции и найдите \(f(-4)\), \(f(0)\), \(f(4)\):
\[
 f(x)=\begin{cases}
 -\dfrac8x,&x<-2\ \text{или}\ x>2,\\[4pt]
 -2x,&-2\le x\le2.
 \end{cases}
\]}

\clearpage
\textbf{Работа с графиком и параметром}
\task{6}{Постройте график функции. Укажите координаты всех граничных закрашенных и выколотых точек, а также значения \(x\), где функция не определена:
\[
 h(x)=\begin{cases}
 -x+2,&x\le-3\ \text{или}\ x\ge3,\\[4pt]
 \dfrac3x,&-3<x<0\ \text{или}\ 0<x<3.
\end{cases}
\]}

\task{7}{Для функции ниже найдите все \(k\), при которых прямая \(y=kx\) имеет с её графиком ровно три общие точки:
\[
 p(x)=\begin{cases}
 -\dfrac{12}{x},&x<-3\ \text{или}\ x>2,\\[4pt]
 -2x-2,&-3\le x\le2.
\end{cases}
\]}

\task{8}{Рассмотрите прямые \(y=3x+4\) и \(y=kx\). При каких \(k\) у них нет общих точек? При каких \(k\) --- одна? Возможны ли две общие точки?}

\task{9}{Для функции из задания 6 найдите все \(k\), при которых прямая \(y=kx\) \emph{не имеет ни одной} общей точки с её графиком.}

\task{10}{Дана функция
\[
 q(x)=\begin{cases}
 -\dfrac4x,&x<-2\ \text{или}\ x>2,\\[4pt]
 -2x,&-2\le x\le2.
\end{cases}
\]
Найдите \(k\), при которых прямая \(y=kx\) имеет с её графиком ровно три общие точки. При каком \(k\) общих точек бесконечно много?}

\clearpage
\secname{Ответы к тренировочному блоку}
\renewcommand{\arraystretch}{1.38}
\begin{center}
\begin{tabularx}{\linewidth}{@{}>{\raggedright\arraybackslash}p{11mm}X@{}}
\toprule
№ & Ответ\\
\midrule
1 & \(x=\dfrac12,\ \dfrac13\).\\
2 & \(x=3,\ -\dfrac32\).\\
3 & \(72\) км.\\
4 & \(60\) км.\\
5 & \(f(-4)=2,\quad f(0)=0,\quad f(4)=-2\). В граничных точках фрагменты соединяются.\\
6 & Закрашены \((-3;5)\), \((3;-1)\); выколоты для гиперболы \((-3;-1)\), \((3;1)\). Функция не определена при \(x=0\).\\
7 & \(\displaystyle k\in\left(-\frac43;0\right)\).\\
8 & Нет точек при \(k=3\); одна при \(k\ne3\); двух быть не может.\\
9 & \(\displaystyle k\in\left(-\infty;-\frac53\right)\cup\{-1\}\cup\left(-\frac13;\frac13\right]\).\\
10 & Ровно три: \(\displaystyle k\in(-1;0)\). Бесконечно много: \(k=-2\).\\
\bottomrule
\end{tabularx}
\end{center}

\vspace{8mm}
\secname{Итоговая памятка}
\renewcommand{\arraystretch}{1.25}
\begin{tabularx}{\linewidth}{@{}>{\raggedright\arraybackslash}p{43mm}X@{}}
\toprule
Сигнал в задаче & Действие\\
\midrule
Повторяется \(1/x\) & Введите \(t=1/x\), решите и вернитесь к \(x\).\\
Есть остановка & Различайте время после выезда и время движения.\\
Функция кусочная & Стройте каждый фрагмент только на его промежутке.\\
Есть строгая граница & Отметьте выколотую точку, если её не включает другой фрагмент.\\
Нужно число пересечений & Решайте \(f_i(x)=kx\), фильтруйте корни по областям.\\
Параметр меняется & Исследуйте критические наклоны и особые случаи.\\
\bottomrule
\end{tabularx}
\vfill
\begin{center}
{\small\color{darkaccent}Лёвин Артём Александрович\quad·\quad эксклюзивно для Софии Кальней}
\end{center}
\end{document}
